\part[3]\hfill

  Here are two implementations of a table lookup, which behave identically when
  the key is present:
  \begin{codeblock}
    exception NotFound
    fun findExn k []            = raise NotFound
      | findExn k ((k',v)::rest) =
          if k = k' then v else findExn k rest

    fun findOpt k []            = NONE
      | findOpt k ((k',v)::rest) =
          if k = k' then SOME v else findOpt k rest
  \end{codeblock}

  Which of these \textbf{forces the caller} to account for the possibility that
  the key is absent? Explain the difference in one or two sentences.

  \begin{solutionorbox}[12em]
    \code{findOpt}. Its result type \code{'a option} makes the possibility of
    failure \textbf{visible in the type}: to use the result, the caller must
    inspect it (case on \code{SOME}/\code{NONE}), so the absent case cannot be
    ignored silently.

    \code{findExn} returns an \code{'a} directly; the possibility of
    \code{NotFound} is invisible in its type, so a caller can use the result as
    if the key were always present --- and the program simply crashes at runtime
    if it is not. The option version encodes the precondition in the type; the
    exception version relies on the caller to remember it.
  \end{solutionorbox}
